\section*{NeurIPS Paper Checklist}

\begin{enumerate}

\item {\bf Claims}
    \item[] Question: Do the main claims made in the abstract and introduction accurately reflect the paper's contributions and scope?
    \item[] Answer: \answerYes{}
    \item[] Justification: The abstract and introduction distinguish the central graded-material study, its reused-identity RGB and comparator analyses, earlier fresh cohorts, and the new 64-training/64-test paired-motion families. Native decoding, unchanged-decoder transfer, and metadata predictability are reported separately. No comparison is presented as identifying a mechanism of persistent object memory.


\item {\bf Limitations}
    \item[] Question: Does the paper discuss the limitations of the work performed by the authors?
    \item[] Answer: \answerYes{}
    \item[] Justification: The interpretation subsection explains finite reset windows, remaining differences in resolution and capacity, synthetic versus physical occlusion, limited model and scene breadth, deterministic motion predictability, central-cohort reuse, and conditional intervals. Historical BGR results are kept distinct from verified RGB results.


\item {\bf Theory assumptions and proofs}
    \item[] Question: For each theoretical result, does the paper provide the full set of assumptions and a complete (and correct) proof?
    \item[] Answer: \answerNA{}
    \item[] Justification: The paper presents an empirical study and does not introduce theoretical results, theorems, or formal proofs.


    \item {\bf Experimental result reproducibility}
    \item[] Question: Does the paper fully disclose all the information needed to reproduce the main experimental results of the paper to the extent that it affects the main claims and/or conclusions of the paper (regardless of whether the code and data are provided or not)?
    \item[] Answer: \answerYes{}
    \item[] Justification: The methods and appendices specify input conventions, checkpoint provenance, alignment filters, exact schedules, family and identity splits, training-only selection, and inferential units. The compact snapshot includes held-out predictions, targets, folds, configurations, hashes, and CPU code that reconstructs the reported statistics without inference. Large activations and model weights are separate inputs for full GPU replication.


\item {\bf Open access to data and code}
    \item[] Question: Does the paper provide open access to the data and code, with sufficient instructions to faithfully reproduce the main experimental results, as described in supplemental material?
    \item[] Answer: \answerNo{}
    \item[] Justification: A local statistical-reproduction snapshot preserves the preceding cohorts and adds transfer, opacity, stronger history comparators, and the paired-motion extension. Anonymous publication of this updated snapshot has not been verified. The existing project link alone does not establish reviewer access to the new artifacts.


\item {\bf Experimental setting/details}
    \item[] Question: Does the paper specify all the training and test details (e.g., data splits, hyperparameters, how they were chosen, type of optimizer) necessary to understand the results?
    \item[] Answer: \answerYes{}
    \item[] Justification: The methodology and appendices specify all cohorts, exact input histories, motion schedules, the fixed representation, grouped splits, ridge selection, baseline information, and source-only transfer. Capture-engineering amendments and superseded artifacts are retained; test prediction outcomes were not used to change the analysis.


\item {\bf Experiment statistical significance}
    \item[] Question: Does the paper report error bars suitably and correctly defined or other appropriate information about the statistical significance of the experiments?
    \item[] Answer: \answerYes{}
    \item[] Justification: Central intervals resample whole identities; new motion intervals resample whole paired-variant families, stratified by layout and motion type. Both use 10,000 draws. The intervals condition on fitted probes, are descriptive, and are not simultaneous bounds across every comparison. The paper distinguishes identity or family averages from pooled frame errors.


\item {\bf Experiments compute resources}
    \item[] Question: For each experiment, does the paper provide sufficient information on the computer resources (type of compute workers, memory, time of execution) needed to reproduce the experiments?
    \item[] Answer: \answerNo{}
    \item[] Justification: The new follow-up reports GPU, CPU, memory, software, measured extraction and fitting times, and preparation time separately. Recovered receipts identify the historical NVIDIA L4 extraction environment. Complete total-project and earlier-development costs are not established, so this answer remains No for the full history rather than implying all historical compute was measured.


\item {\bf Code of ethics}
    \item[] Question: Does the research conducted in the paper conform, in every respect, with the NeurIPS Code of Ethics \url{https://neurips.cc/public/EthicsGuidelines}?
    \item[] Answer: \answerYes{}
    \item[] Justification: The study uses a controlled simulated environment and does not involve human subjects, personal data, or scraped content. The research and anonymous submission follow the NeurIPS Code of Ethics.


\item {\bf Broader impacts}
    \item[] Question: Does the paper discuss both potential positive societal impacts and negative societal impacts of the work performed?
    \item[] Answer: \answerNA{}
    \item[] Justification: This is a diagnostic study in a simulated environment, not a deployed decision system. Its immediate benefit is clearer evaluation of world-model readouts; extrapolation to reliable physical decision-making is explicitly outside the measured claims.


\item {\bf Safeguards}
    \item[] Question: Does the paper describe safeguards that have been put in place for responsible release of data or models that have a high risk for misuse (e.g., pre-trained language models, image generators, or scraped datasets)?
    \item[] Answer: \answerNA{}
    \item[] Justification: The released benchmark and evaluation materials do not introduce a new high-risk generative model or contain personal or human-subject data.


\item {\bf Licenses for existing assets}
    \item[] Question: Are the creators or original owners of assets (e.g., code, data, models), used in the paper, properly credited and are the license and terms of use explicitly mentioned and properly respected?
    \item[] Answer: \answerYes{}
    \item[] Justification: MineWorld and Fabric are credited under their respective MIT and Apache-2.0 licenses. VideoMAE-S and its official pretrained source are credited under CC-BY-NC-4.0 and used for noncommercial research. Minecraft game files and pretrained model weights are not redistributed in the compact snapshot. Asset provenance and upstream links are included.


\newpage
\item {\bf New assets}
    \item[] Question: Are new assets introduced in the paper well documented and is the documentation provided alongside the assets?
    \item[] Answer: \answerYes{}
    \item[] Justification: The snapshot documents schemas, exact schedules and splits, preprocessing, model provenance, prediction files, uncertainty calculations, validation, and execution commands. It distinguishes CPU prediction-to-statistics reconstruction from full inference and from the separately pending verification of anonymous distribution.


\item {\bf Crowdsourcing and research with human subjects}
    \item[] Question: For crowdsourcing experiments and research with human subjects, does the paper include the full text of instructions given to participants and screenshots, if applicable, as well as details about compensation (if any)? 
    \item[] Answer: \answerNA{}
    \item[] Justification: The study does not involve crowdsourcing, human participants, or data collected from human subjects.


\item {\bf Institutional review board (IRB) approvals or equivalent for research with human subjects}
    \item[] Question: Does the paper describe potential risks incurred by study participants, whether such risks were disclosed to the subjects, and whether Institutional Review Board (IRB) approvals (or an equivalent approval/review based on the requirements of your country or institution) were obtained?
    \item[] Answer: \answerNA{}
    \item[] Justification: The study does not involve human participants or human-subject data, so IRB approval or an equivalent review was not required.


\item {\bf Declaration of LLM usage}
    \item[] Question: Does the paper describe the usage of LLMs if it is an important, original, or non-standard component of the core methods in this research? Note that if the LLM is used only for writing, editing, or formatting purposes and does \emph{not} impact the core methodology, scientific rigor, or originality of the research, declaration is not required.
    \item[] Answer: \answerYes{}
    \item[] Justification: LLM assistance supported the prospective diagnostic design, implementation, analysis scripts, debugging, and manuscript preparation, as disclosed in the diagnostic appendix. It did not generate the capture images, targets, or reported numerical findings. Measurements are bound to executed captures and model predictions with explicit validation checks.

\end{enumerate}
